\chapter{Conclusion}
The quantum state vector is a universal coordinate system, not always the computationally intrinsic one. Conserved charges, finite-group representations, permutation orbits, few-particle sectors, and operator packets can expose much smaller dynamics. When the packet is efficiently discoverable, polynomial in size, closed under the gates, and compatible with the measurement, the quantum computation collapses exactly to a deterministic classical polynomial-time simulation.

The strongest current examples are concrete. Fixed charge at constant particle number is polynomial. Exact charge conservation is not necessary: moving packet schedules, efficiently computable moving frames, and bounded unions of low-excitation sectors also collapse to polynomial classical simulation. Fully symmetric fixed-local-dimension dynamics is polynomial. A stronger mixed-state theorem compresses all permutation-invariant Schur sectors into a polynomial commutant record, with $O(n^3)$ density-matrix parameters for qubits under efficient block access. Polynomial-orbit walks are polynomial under quotient-access hypotheses. Few-particle bosonic and fermionic systems are polynomial under efficient updates, while growing particle number gives a sharp progression through quasipolynomial, subexponential, and exponential sector sizes. Stabilizer and Gaussian representations show that an exponential packet can also be easy when its internal states admit compact closed coordinates. Recorded number-conserving benchmarks show large memory and repeated-application gains with zero reconstruction error at the recorded precision.

The boundaries are equally sharp. Parity, modulo-four charge, and half filling do not generally change exponential scaling. Entanglement can remain hard inside one packet. A syndrome label does not compress logical information. A block encoding does not provide a classical access model. Global Fourier interference may be fully structured yet resist small packetization. Symmetry collapse is therefore not a slogan that all quantum computation is classical. It is a theorem architecture for identifying substantial subclasses that are.

The core research invariant is the complete classical representation record
\[
\bigl(
\text{discovery cost},
N_{\rm packets},
D_{\rm internal},
S_{\rm transfer},
L_{\rm leakage},
\text{measurement cost}
\bigr).
\]
A future compiler that computes and certifies this record can decide when the full state vector is unnecessary, when a hybrid representation is preferable, and when symmetry alone is insufficient.
